\chapter{Overactive thyroid}
\index{Overactive thyroid}

\medskip
\noindent\textit{Educational information; seek professional advice for individual care.}

\section*{What it is}
An overactive thyroid, or hyperthyroidism, occurs when the thyroid produces more thyroid hormone than the body needs. Causes include Graves disease, overactive thyroid nodules, thyroid inflammation, excess iodine, or too much thyroid hormone medicine.

\section*{Symptoms and complications}
Possible symptoms include weight loss despite appetite, fast or irregular heartbeat, tremor, sweating, heat intolerance, anxiety, poor sleep, tiredness, muscle weakness, frequent bowel movements, neck swelling, or menstrual and fertility changes. Untreated disease can affect the heart, bones, muscles, eyes, and pregnancy.

\section*{Diagnosis}
Symptoms alone cannot confirm the diagnosis. Assessment usually includes thyroid blood tests and may include antibody tests, ultrasound, or a thyroid uptake scan to identify the cause. Tell the clinician about pregnancy, heart disease, iodine-containing medicines or supplements, and all thyroid medicines.

\section*{Treatment}
Treatment depends on the cause and may include beta-blockers for selected symptoms, antithyroid medicines, radioactive iodine, or surgery. These treatments have different risks and are not interchangeable. Follow-up blood tests are needed, and treatment can sometimes lead to an underactive thyroid.

\section*{When to seek urgent care}
Seek emergency help for chest pain, severe breathlessness, fainting, confusion, a very fast or irregular heartbeat, high fever with marked agitation, or rapidly worsening illness. If taking an antithyroid medicine, promptly report fever, persistent sore throat, jaundice, severe weakness, or unusual bruising.

\section*{Sources}
\begin{itemize}
\item National Institute of Diabetes and Digestive and Kidney Diseases: \url{https://www.niddk.nih.gov/health-information/endocrine-diseases/hyperthyroidism}
\end{itemize}
